% Part: first-order-logic
% Chapter: introduction
% Section: satisfaction

\documentclass[../../../include/open-logic-section]{subfiles}

\begin{document}

\olfileid{fol}{int}{sat}

\olsection{पूर्ति}

!!{formula}s म्हणजे काय हे माहीत झाल्यावर आपण लगेच प्रथम-क्रम
तर्कशास्त्राच्या चिन्हार्थमीमांसेकडे वळू शकतो: येथे मूलभूत व्याख्या
!!a{structure} ची आहे. आपल्या साध्या भाषेसाठी !!a{structure}~$\Struct M$
चे फक्त तीन घटक असतात: \emph{!!{domain}} म्हणतात तो अरिक्त संच
$\Domain{M}$, $\Struct M$ मध्ये $\Obj a$ ने निर्देशित केलेली गोष्ट, आणि
$\Struct M$ मध्ये $\Obj P$ ज्या गोष्टींबाबत सत्य आहे त्या गोष्टी.
$\Obj a$ ने निर्देशित केलेली वस्तू~$\Assign{\Obj a}{M}$ ने, तर
$\Obj P$ ज्या गोष्टींबाबत सत्य आहे त्यांचा संच~$\Assign{\Obj P}{M}$ ने
दर्शवतात. !!^a{structure}~$\Struct{M}$ मध्ये याच तीन गोष्टी असतात:
$\Domain{M}$, $\Assign{\Obj a}{M} \in \Domain{M}$ आणि
$\Assign{\Obj P}{M} \subseteq \Domain{M}$. सर्वसाधारण स्थिती अधिक
गुंतागुंतीची असेल, कारण त्यात अनेक !!{predicate}s आणि !!{constant}s
असतील, !!{predicate}s ना एकाहून अधिक स्थाने असू शकतील, आणि
!!{function}s ही असतील.
%READERNOTE{OLFOL-002: गोठवलेल्या इंग्रजी मजकुरात अनेक स्थाने असू शकणारी चिन्हे constants म्हटली आहेत. स्थिरचिन्हे एकच वस्तू निर्देशित करतात; अनेक स्थाने असू शकणारी चिन्हे predicates असल्यामुळे मराठी मजकुरात ते दुरुस्त केले आहे.}

ज्या !!{formula}s मध्ये !!{variable}s नाहीत त्यांच्यासाठी पूर्तिची
व्याख्या देण्यास हे पुरेसे आहे. आपण !!{formula}s ज्या रीतीने परिभाषित
केली तिचे प्रतिबिंब दाखवणारी विगमनाधारित व्याख्या देण्याची कल्पना आहे.
एखाद्या आण्विक सूत्राची~$\Struct{M}$ मध्ये पूर्ति केव्हा होते हे आपण
सांगतो आणि मग, उदाहरणार्थ, $!A$ ची~$\Struct{M}$ मध्ये पूर्ति होते की
नाही याच्या आधारे $\lnot !A$ ची~$\Struct{M}$ मध्ये पूर्ति केव्हा होते
हे सांगतो. उदाहरणार्थ, पुढीलप्रमाणे व्याख्या करता येईल:
\begin{enumerate}
  \item $\Atom{\Obj P}{\Obj a}$ ची~$\Struct{M}$ मध्ये पूर्ति तेव्हा आणि
  केवळ तेव्हाच होते, जेव्हा $\Assign{\Obj a}{M} \in \Assign{\Obj P}{M}$.
  \item $\lnot !A$ ची~$\Struct{M}$ मध्ये पूर्ति तेव्हा आणि केवळ तेव्हाच
  होते, जेव्हा $!A$ ची~$\Struct{M}$ मध्ये पूर्ति होत नाही.
  \item $(!A \land !B)$ ची~$\Struct{M}$ मध्ये पूर्ति तेव्हा आणि केवळ
  तेव्हाच होते, जेव्हा $!A$ ची~$\Struct{M}$ मध्ये पूर्ति होते आणि
  $!B$ ची~$\Struct{M}$ मध्येही पूर्ति होते.
\end{enumerate}
$\Domain{M} = \{0, 1, 2\}$, $\Assign{\Obj a}{M} = 1$ आणि
$\Assign{\Obj P}{M} = \{1, 2\}$ असे समजू. ही व्याख्या आपल्याला
$\Atom{\Obj P}{\Obj a}$ ची~$\Struct{M}$ मध्ये पूर्ति होते असे सांगेल
(कारण $\Assign{\Obj a}{M} =  1 \in \{1,2\} = \Assign{\Obj P}{M}$).
ती पुढे असे सांगते की $\lnot \Atom{\Obj P}{\Obj a}$ ची~$\Struct{M}$
मध्ये पूर्ति होत नाही; म्हणून $\lnot\lnot \Atom{\Obj P}{\Obj a}$ ची
पूर्ति होते आणि $(\lnot \Atom{\Obj P}{\Obj a} \land \Atom{\Obj P}{\Obj a})$
ची पूर्ति होत नाही; आणि असेच पुढे जाता येते.

संख्यापकांसाठी व्याख्या द्यायची झाल्यावर अडचण निर्माण होते: आपल्याला
``$\lexists[\Obj v_0][\Atom{\Obj P}{\Obj v_0}]$ ची पूर्ति तेव्हा आणि
केवळ तेव्हाच होते, जेव्हा $\Atom{\Obj P}{\Obj v_0}$ ची पूर्ति होते''
असे काहीतरी म्हणायचे आहे. पण !!{structure}~$\Struct{M}$ हे
!!{variable}s विषयी काय करायचे ते सांगत नाही. आपल्याला प्रत्यक्षात
$\Atom{\Obj P}{\Obj v_0}$ ची \emph{$\Obj v_0$ च्या एखाद्या मूल्यासाठी}
पूर्ति होते असे म्हणायचे आहे. हे काटेकोर करण्यासाठी फक्त $\Obj a$ लाच
नव्हे, तर~$\Obj v_0$ लाही~$\Domain{M}$ मधील !!{element}s नेमून देण्याचा
मार्ग हवा. यासाठी आपण !!{variable} ची \emph{मूल्यांकने} आणतो.
!!^a{variable} चे मूल्यांकन म्हणजे !!{variable}s ना~$\Domain{M}$ मधील
!!{element}s कडे पाठवणारे फलन~$s$ (आपल्या उदाहरणात, $0$, $1$ किंवा~$2$
यांपैकी एका घटकाकडे). एखाद्या !!a{formula} मध्ये कोणती !!{variable}s
येतील हे आपल्याला आधी माहीत नसल्यामुळे, $s$ कोणत्या !!{variable}s ना
मूल्ये नेमून देते यावर मर्यादा घालता येत नाही. सोपा उपाय म्हणजे $s$ ने
\emph{सर्व} !!{variable}s~$\Obj v_0$, $\Obj v_1$, \dots\@ यांना मूल्ये
नेमून द्यावीत अशी अट ठेवणे. आपल्याला लागतील तीच मूल्ये आपण वापरू.
%READERNOTE{OLFOL-003: याच उदाहरणात प्रांत {0, 1, 2} असा ठरवला आहे; गोठवलेल्या इंग्रजी मजकुरातील {1, 2, 3} ही यादी 0 वगळते आणि प्रांताबाहेरील 3 समाविष्ट करते. मराठी मजकुरात प्रांताशी जुळणारी {0, 1, 2} ही यादी वापरली आहे.}

!!{formula}s ची पूर्ति फक्त !!a{structure} च्या सापेक्ष परिभाषित
करण्याऐवजी, आपण ती !!a{structure}~$\Struct{M}$ \emph{आणि}
!!a{variable} चे मूल्यांकन~$s$ या दोन्हींच्या सापेक्ष परिभाषित करू आणि
संक्षेपाने $\Sat{M}{!A}[s]$ असे लिहू. !!{variable}s असलेली आण्विक
!!{formula}s हाताळण्यासाठी आपल्या व्याख्येत आता एक अतिरिक्त उपवाक्य
असेल:
\begin{enumerate}
  \item $\Sat{M}{\Atom{\Obj P}{\Obj a}}[s]$ तेव्हा आणि केवळ तेव्हाच,
  जेव्हा $\Assign{\Obj a}{M} \in \Assign{\Obj P}{M}$.
  \item $\Sat{M}{\Atom{\Obj P}{\Obj v_i}}[s]$ तेव्हा आणि केवळ तेव्हाच,
  जेव्हा $s(\Obj v_i) \in \Assign{\Obj P}{M}$.
  \item $\Sat{M}{\lnot !A}[s]$ तेव्हा आणि केवळ तेव्हाच, जेव्हा
  $\Sat{M}{!A}[s]$ नाही.
  \item $\Sat{M}{(!A \land !B)}[s]$ तेव्हा आणि केवळ तेव्हाच, जेव्हा
  $\Sat{M}{!A}[s]$ आणि $\Sat{M}{!B}[s]$ दोन्ही आहेत.
\end{enumerate}
ठीक आहे, यामुळे एक अडचण सुटते: $s(\Obj v_0)$ या मूल्यासाठी
$\Struct{M}$ हे $\Atom{\Obj P}{\Obj v_0}$ ची पूर्ति केव्हा करते हे आता
सांगता येते. $\lexists[\Obj v_0][\Atom{\Obj P}{\Obj v_0}]$ साठी योग्य
व्याख्या मिळवण्यासाठी आणखी एक गोष्ट करावी लागेल:
$\Sat{M}{\lexists[\Obj v_0][\Atom{\Obj P}{\Obj v_0}]}[s]$ तेव्हा आणि
केवळ तेव्हाच असावे, जेव्हा~$\Obj v_0$ ला मूल्य नेमून देणाऱ्या
\emph{एखाद्या} पद्धती~$s'$ साठी
$\Sat{M}{\Atom{\Obj P}{\Obj v_0}}[s']$ असते. पण~$\Obj v_0$ ला नेमलेले
मूल्य $s(\Obj v_0)$ ने निर्देशित केलेले मूल्यच असणे आवश्यक नाही.
त्यासाठी आपण एक संकेतपद्धती आणू: $m \in \Domain{M}$ असेल, तर
$\Subst{s}{m}{\Obj v_0}$ हे~$s$ सारखेच असणारे (फक्त~$\Obj v_0$ सोडून
इतर सर्व !!{variable}s साठी) पण~$\Obj v_0$ ला~$m$ नेमून देणारे
मूल्यांकन मानू. आता आपली व्याख्या अशी देता येईल:
\begin{enumerate}\setcounter{enumi}{4}
  \item $\Sat{M}{\lexists[\Obj v_i][!A]}[s]$ तेव्हा आणि केवळ तेव्हाच,
  जेव्हा एखाद्या~$m \in \Domain{M}$ साठी
  $\Sat{M}{!A}[\Subst{s}{m}{\Obj v_i}]$.
\end{enumerate}
हे योग्य रीतीने चालते का? प्रत्येक~$i \in \Nat$ साठी
$s(\Obj v_i) = 0$ ठेवू या. एखादा $m \in \Domain{M}$ असा असेल की
$\Sat{M}{\Atom{\Obj P}{\Obj v_0}}[\Subst{s}{m}{\Obj v_0}]$, तेव्हा आणि
केवळ तेव्हाच $\Sat{M}{\lexists[\Obj v_0][\Atom{\Obj P}{\Obj v_0}]}[s]$.
आणि असा घटक आहे: आपण $m = 1$ किंवा $m = 2$ निवडू शकतो. $s$ नेच
~$\Obj v_0$ ला नेमून दिलेले $s(\Obj v_0)$ हे मूल्य---या उदाहरणात,
$0$---हे काम करत नसले तरी हे सत्य आहे, हे लक्षात घ्या. आपल्याकडे
$\Sat{M}{\Atom{\Obj P}{\Obj v_0}}[\Subst{s}{1}{\Obj v_0}]$ आहे, पण
$\Sat{M}{\Atom{\Obj P}{\Obj v_0}}[s]$ नाही.

हे गोंधळात टाकणारे आणि अवजड वाटत असेल, तर ते तसे आहे. पण सर्व
!!{formula}s साठी पूर्तिची काटेकोर, विगमनाधारित व्याख्या देण्यासाठी ही
अतिरिक्त गुंतागुंत आवश्यक आहे आणि चिन्हार्थविषयक संकल्पना काटेकोरपणे
परिभाषित करण्यासाठी आपल्याला अशाच काहीतरी पद्धतीची गरज आहे. हे
करण्याचे इतरही मार्ग आहेत, पण ते सर्व तितकेच (अन)सुबक आहेत.

\end{document}
